\section{Conclusion}
\label{sec:conclusion}

This thesis set out to design, implement, and validate a high-performance
simulation sandbox for developing and testing decentralised control algorithms
for autonomous marine drone swarms. The work was motivated by the absence
of any existing platform that simultaneously addresses realistic marine
dynamics, swarm-scale performance, physics-based communication constraints,
and a direct path to physical hardware deployment. Major contribution of this project
is the deep mathematical foundation underlying the simulator. Every simulated system and component is described by a precise mathematical model. In other projects the models of sensors, actuators etc, were usualy igonored. In this project, the mathematical models of the drone ship physics, wave dynamics, LoRa radio propagation, and AIS signal characteristics were derived from first principles.

The simulator was validated through a series of experiments designed to answer three core research questions, which are revisited in the following sections. The results confirm that the platform meets its design goals and provides a practical tool for developing and testing swarm control algorithms under realistic conditions.


\subsection*{RQ1 — Simulator Fidelity and Scalability}

The Unity3D-based simulator successfully models the key physical phenomena
relevant to autonomous USV operations. The drone ship physics model
combines a multi-point buoyancy system, wave-slope dynamics derived from a
Pierson–Moskowitz spectral wave model, a keel-drag curve, propeller slipstream
effects, and a full hydrodynamic rudder model. Virtual sensor outputs — GPS,
IMU, magnetometer, AIS transponder and receiver, LoRa and time-of-flight
ranging — are augmented with hardware-realistic noise models, ensuring that
control algorithms are tested against the same signal degradation they would
encounter at sea.

Scalability benchmarks, conducted on a mid-range laptop with integrated
graphics (AMD Ryzen\texttrademark~5 7530U, Radeon\texttrademark~Graphics),
demonstrated that the simulator sustains real-time performance at 60~FPS with
up to 71 simulated vessels, and remains above the minimum acceptable
threshold of 30~FPS with up to 103 vessels running concurrently inside the
Unity Editor. A compiled standalone build is expected to push this limit
further. The communication architecture, based on lightweight UDP/JSON
packets and a dedicated port pair per drone agent, introduces negligible
overhead and allows an arbitrary number of independent controller instances to
operate concurrently without port conflicts. Transitioning from
Software-in-the-Loop to Hardware-in-the-Loop requires modifications to only
two configuration files, confirming that the platform provides a practical path
to physical deployment.

\subsection*{RQ2 — Sensor Fusion Resilience}

A Kalman Filter fusing GPS, IMU, compass heading, and LoRa time-of-flight
ranging was implemented and validated under three conditions. Under normal
operation the filter reduced position RMSE from approximately 1.8~m (GPS
alone) to 1.08~m. When the GPS signal of the tested drone was spoofed with
Gaussian noise of standard deviation 70~m, the filter maintained an RMSE of
approximately 19~m — a more than threefold improvement over the corrupted
GPS reading. In the most challenging scenario, with half the swarm
simultaneously spoofed, the filter achieved an RMSE of approximately
1.6~m by leveraging LoRa inter-drone ranging from the unaffected neighbours,
which is still below the accuracy of the raw GPS in normal conditions.

These results confirm that the EKF-based localisation system provides
acceptable positioning accuracy for swarm coordination even under active GPS
spoofing, and that the cooperative nature of the swarm provides a natural
resilience mechanism that a single-vehicle system cannot replicate.

\subsection*{RQ3 — Swarm Search Performance}

A decentralised grid-based search algorithm, coordinated exclusively over
LoRa, was validated in a realistic maritime scenario: a $3\,\text{km} \times
3\,\text{km}$ anchorage zone in the Gulf of Gdańsk, discretised into
$300\,\text{m} \times 300\,\text{m}$ cells, with five vessels moving according
to replayed real AIS traffic. A swarm of 18 drones achieved 100\,\% area
coverage in approximately 45 minutes. A reduced swarm of 9 drones completed
the same task in approximately 53 minutes — only 8 minutes longer — which
suggests that at this scale the collision avoidance overhead and cell density
are the binding constraints rather than drone count alone. In both cases no
inter-drone collisions were recorded, and all AIS vessel avoidance manoeuvres
were executed successfully without disrupting overall coverage progress.

The decentralised cell-assignment scoring function, which combines staleness,
a contention penalty propagated via LoRa beacons, and a distance
tie-breaker, proved effective at coordinating the swarm without any central
controller. The LoRa beacon protocol, encoded in at most 36~bytes per
transmission, remained within the duty-cycle limits of the Semtech SX1280 at
the spreading factors typical of maritime deployments.

\subsection*{Summary of Contributions}

The primary contributions of this work are:

\begin{enumerate}
    \item A Unity3D simulation environment for marine USV
    swarms, capable of real-time simulation of over 100 vessels on commodity
    hardware, with no dependency on ROS or a specific operating system.

    \item A physics model that captures the essential hydrodynamics of a
    small surface vessel, including buoyancy, wave-slope dynamics, keel
    resistance, and rudder lift and drag, implemented efficiently enough to
    run at 200~Hz inside the physics integration loop.

    \item A hardware-realistic LoRa radio propagation model, including
    free-space path loss, additive noise, half-duplex enforcement, collision
    modelling, and time-of-flight ranging.

    \item An AIS data replay pipeline that ingests CSV-formatted real-world
    AIS recordings, reconstructs smooth vessel trajectories, and injects replayed traffic into the simulator's internal AIS network so that it is indistinguishable from actively simulated vessels.

    \item A personal AIS receiver station based on the EM-TRAK B100 and an
    Orange Pi PC, which produced the real-world dataset used for the
    validation experiments.

    \item A Kalman Filter localisation system that fuses GPS, IMU, compass,
    and LoRa inter-drone range measurements, demonstrated to maintain
    sub-2\,m RMSE under normal conditions and to resist GPS spoofing
    attacks by cooperative ranging.

    \item A decentralised swarm controller comprising a PID course-and-speed
    regulator, swarm anti-collision avoidance system, AIS vessel avoidance system,
    and a staleness-driven grid coverage algorithm coordinated over LoRa.
\end{enumerate}

\subsection*{Limitations and Future Work}

Several limitations remain. The LoRa propagation model assumes a spatially
uniform path-loss exponent and does not account for Doppler shift, multipath
from dynamic wave surfaces, or interference from co-located 2.4~GHz systems
such as Wi-Fi. Also signal in real life could be blocked bt big ship.
The wave model, while spectrally motivated, uses a small fixed
set of harmonics and does not yet couple wind field dynamics to wave growth.
The drone physics model does not incorporate a hull-shape-dependent added
mass tensor, which would be necessary for accurate simulation of fast
manoeuvres. The swarm controller relies entirely on AIS for third-party vessel
detection; vessels that are dark or whose AIS is spoofed are invisible to the
avoidance system.

Future work should address these limitations and extend the platform in
several directions. The most immediate next step is a Hardware-in-the-Loop
validation using the physical Argo-style hardware described in Section~0.4, to
confirm that the SITL-to-HITL migration path works as designed under real
network latency and microcontroller constraints. On the algorithmic side, the
deterministic swarm controller could be replaced or augmented by a
neural-network or genetic-algorithm policy trained using the multi-instance
parallelism that the simulator already supports. Adding a radar or camera
sensor model would remove the dependency on AIS for obstacle detection,
significantly improving safety in cluttered or adversarial environments.
Finally, extending the AIS data collection station to produce longer and
geographically richer datasets would enable more demanding validation
scenarios and facilitate the study of swarm performance as a function of
traffic density.